\section{Interpretation, scale, and limitations}
\label{sec:interpretation}

Equation~\eqref{eq:headline} is a passive spectral-area statement. Resonances can redistribute gravitational response in frequency, internal hybridization can create bright and dark combinations, and matching choices can determine how much of a particular resonance is accessible from a selected local port. What they cannot do, within the retained modal class, is create additional cumulative gravitational coupling trace beyond the endpoint second-moment resource. Free-space propagation then imposes the separate compact TT channel ceiling.

The minimum in Eqs.~\eqref{eq:asymptotic-headline} and \eqref{eq:headline} is physically important. Either endpoint is a cut: a large source resource cannot compensate for a receiver whose total retained gravitational coupling trace is smaller, and vice versa. Elastic constants, quality factors, and internal modal structure can change where the response appears in frequency, but after optimization they do not survive in the leading cumulative ceiling; only the endpoint mass distributions through $I_2$ and the separated TT channel remain.

\subsection{Scale and a resonant-bar tightness check}

For a uniform sphere of mass $M$ and radius $a$,
\begin{equation}
I_2=\frac35Ma^2,
\end{equation}
so Eq.~\eqref{eq:headline} becomes, for equal endpoints,
\begin{equation}
\Gamma_{\rm coh}
\lesssim
\frac{5G\omega_0^2Ma^2}{4c^3R^2}.
\label{eq:spherebound}
\end{equation}
As an intentionally generous macroscopic illustration, take $M=10^3\,\mathrm{kg}$, $a=1\,\mathrm{m}$, and $f_0=\omega_0/(2\pi)=1\,\mathrm{kHz}$. The compactness condition is then extremely well satisfied, $k_0a\simeq2.1\times10^{-5}$. Choosing $k_0R=100$ requires $R\simeq4.77\times10^6\,\mathrm{m}$ and gives
\begin{equation}
\Gamma_{\rm coh}\lesssim5.4\times10^{-39}\,\mathrm{s}^{-1}.
\end{equation}
Equation~\eqref{eq:finiteTT} shows that at $k_0R=100$ the finite-distance correction to the compact TT propagation power coefficient is about $0.020\%$; the example therefore does not rely on an unquantified propagation asymptotic. The separate compact-source and narrowband approximations remain part of the stated model.

A simple bar calculation also tests whether the endpoint $20/3$ closure is vacuous by many orders of magnitude. Consider an ideal slender free--free uniform bar of length $L$ and mass $M$, neglecting transverse dimensions, with fundamental longitudinal displacement $w_x(x)=\sin(\pi x/L)$ for $-L/2\le x\le L/2$. Then
\begin{equation}
\mu=\frac{M}{2},\qquad I_2=\frac{ML^2}{12},
\end{equation}
and direct evaluation of Eq.~\eqref{eq:qn} gives
\begin{equation}
\frac{q:q}{\mu}=\frac{64ML^2}{3\pi^4}.
\end{equation}
Consequently this single idealized mode occupies the fraction
\begin{equation}
\boxed{
\frac{(q:q)/\mu}{(20/3)I_2}
=\frac{192}{5\pi^4}
\simeq0.394
}
\label{eq:barfraction}
\end{equation}
of the complete unweighted endpoint sum-rule ceiling. Its axisymmetric quadrupole has maximum directivity $15/8$, below the unrestricted STF ceiling $5/2$. Thus the endpoint and angular constants are not loose by dozens of orders of magnitude even for the canonical resonant-bar geometry. This calculation does \emph{not} imply that an existing resonant-mass detector approaches the complete two-ended $\Gamma_{\rm coh}$ ceiling: local-port accessibility, nongravitational losses, matching, omitted modal sectors, and the second endpoint can all reduce the realized transfer. The bound is therefore not claimed to be the active engineering constraint for present detectors. Classic cylindrical-bar cross-section theory provides the relevant historical comparison class \cite{PaikWagoner1976}.

The sphere number is not a waiting time for an event or a bit: $\Gamma_{\rm coh}$ is a spectral area. It makes the scale of the passive far-zone ceiling concrete. In ordinary macroscopic mechanical regimes, the simultaneous compactness and wave-zone assumptions place the theorem in an extraordinarily weak-coupling domain. Its value is therefore primarily structural: it identifies what passive resonance complexity cannot change once the model's endpoint resources and propagation class have been fixed.

\subsection{Relation to prior work and scope}

The ingredients have substantial precedents. Resonant-mass absorption cross sections and integrated response are historical \cite{PaikWagoner1976,Aguiar2011}; arbitrary-body multimode gravitational response is established \cite{Lobo1995}; material-response dispersion and sum-rule methods have been developed \cite{Srivastava2003}; and complete gravitational generator--receiver analyses exist \cite{Rudenko2003}. Generic wave-channel decompositions, response-plus-propagation bounds, passive $H_2$ methods, and multiple-scattering composition are likewise established \cite{Miller2000,Molesky2020,BarasBrockett1975,Opmeer2013,Redheffer1962}.

The closest classical engineering analogy is gain--bandwidth theory, but the correspondence is not identity. Fano's broadband-matching limits constrain the achievable reflection coefficient of a prescribed load through logarithmic integral relations \cite{Fano1950}; Chu and Harrington constrain antenna stored energy, $Q$, directivity, gain, efficiency, and electrical size \cite{Chu1948,Harrington1960}. Those results establish that passive matching or superdirectivity cannot evade generic bandwidth and size tradeoffs. The passive $H_2$ cut in Sec.~\ref{sec:passive} belongs to that broad family of ideas and is not the gravity-specific contribution. What is added here is the closure of each retained gravitational port trace to the same mechanical scalar $I_2$, followed by the compact TT propagation coefficient and the two-ended minimum over source and receiver resources. The historical resonant-mass integrated cross sections are not identified numerically with $\Gamma_{\rm coh}$ because they use different response normalizations.

Recent gravity-communication work addresses different performance criteria and often different interaction regimes: Newtonian communication/noise constraints \cite{KafriMilburnTaylor2015}, LOCC simulation bounds \cite{LamiPedernalesPlenio2024}, state-transfer benchmarks for gravitationally interacting oscillators \cite{ToccaceloAndersenBrask2025}, and an optomechanical gravity-induced thermal-attenuator channel \cite{MariZippilliVitali2026}. These works motivate the broader question of gravity as a mediator, but their quantum-information criteria are not used in the proof. The present quantity is a classical coherent-transfer spectral area, not a channel capacity, entanglement rate, or usability criterion. No equivalent complete passive wave-zone two-ended inertia closure was found in the inspected literature; this is a statement about the search performed, not a priority proof.

The retained-sector restriction is consequential. A finite elastic body has a countable spectrum in the continuum model, while microscopic solids also contain increasingly high-frequency lattice modes beyond that approximation. The present theorem does not prove that high-frequency off-resonant tails are negligible in a chosen operating band. That is a material- and device-specific question. If such a tail cannot be bounded separately, Eq.~\eqref{eq:headline} applies only to the declared retained carrier-scale sector and not to the complete physical endpoint.

The theorem deliberately excludes active gain, inversion, parametric pumping, externally powered feedback, extended phased apertures, added gravitational relays or external cavities, higher-multipole-dominated operation, relativistic or nonlinear matter, arbitrary unbounded PDE boundary ports without admissibility analysis, genuinely non-Markov continua, curved-background focusing, and uncontrolled higher-frequency modal sectors whose off-resonant tails enter the operating band. It is also a compact narrowband carrier-envelope result: $k_0a_{A,B}\ll1$, $k_0R\gg1$, and $B/\omega_0\ll1$. In the near field, $k_0R\lesssim1$, reactive and nonradiative Green-function terms change the distance and frequency structure, so Eq.~\eqref{eq:headline} should not be extrapolated there. A broad absolute-frequency theorem would likewise have to retain the explicit frequency dependence of the matter and propagation resources inside the integral.

Within these limits, the result answers a specific structural question: passive resonant complexity can reshape the gravitational transfer spectrum, but it cannot raise the retained far-zone spectral area above the smaller endpoint inertia resource multiplied by the TT propagation ceiling. Whether that ceiling is the active limitation for any particular device requires a separate device-level coupling and loss model.

\section{Conclusion}

For a separated compact passive gravitational link, passivity reduces the frequency-integrated coherent spectral area to the smaller gravitational coupling trace of the two endpoints. Mass-weighted quadrupole completeness bounds each retained trace by $(4G/3c^5)I_2\Omega^4$, while normalized compact TT propagation supplies the asymptotic coefficient $25/[16(k_0R)^2]$ and an explicit finite-$kR$ correction within the outgoing compact angular model. Their combination gives the precise asymptotic statement Eq.~\eqref{eq:asymptotic-headline}, whose carrier-scale form is
\begin{equation}
\Gamma_{\rm coh}
\lesssim
\frac{25G\omega_0^2}{12c^3R^2}
\min(I_{2,A},I_{2,B}).
\end{equation}
The same endpoint cut extends to countably infinite bounded-port Markov modal sectors satisfying the retained-sector assumptions, and same-two-endpoint passive recurrence alters only subleading wave-zone orders. The result is deliberately narrower than a general statement about gravitational communication or arbitrary material spectra: within its compact, narrowband, passive far-zone class, adding resonances, increasing passive $Q$, or hybridizing modes can redistribute response but cannot manufacture additional cumulative two-ended gravitational spectral area beyond the endpoint mass-distribution resource.
